\documentclass[a4paper,11pt]{report}
\usepackage[utf8]{inputenc}
\usepackage[T1]{fontenc}
\usepackage{lmodern}
\usepackage[ngerman]{babel}
\usepackage[margin=20mm, left=20mm, right=10mm, headheight=15pt, includeheadfoot]{geometry}
\usepackage{fancyhdr}
\usepackage{opensans}
\usepackage{titlesec}
\usepackage{tocloft}
\usepackage{titling}
\usepackage{hyperref}
\usepackage{graphicx}
\usepackage{enumerate}
\usepackage{float}
\usepackage{caption}
\usepackage{listings}
\usepackage{minted}

\usemintedstyle{vs}

\lsbset{
    frame=tb,
    aboveskip=3mm,
    belowskip=3mm,
    columns=flexible,
    basicstyle={\small\ttfamily},
    numbers=none,
    breaklines=true,
    tabsize=3
}

% Set default font to OpenSans
\renewcommand*\familydefault{\sfdefault}

% Page style for cover page
\fancypagestyle{cover}{
    \fancyhf{}
    \renewcommand{\headrulewidth}{0pt}
    \renewcommand{\footrulewidth}{0pt}
}

% Page style for main content
\fancypagestyle{main}{
    \fancyhf{}
    \fancyhead[L]{\theauthor}
    \fancyhead[R]{\subject}
    \fancyfoot[C]{\thepage}
    \renewcommand{\headrulewidth}{0.4pt}
    \renewcommand{\footrulewidth}{0pt}
}

\titleformat{\chapter}[block]
  {\normalfont\Large\bfseries} % change \Large to \large or any other size that fits
  {\thechapter}
  {1em}
  {}

  \titlespacing*{\chapter}{0pt}{*4}{*2.5}

% Cover page
\newcommand{\coverpage}{
    \thispagestyle{cover}
    \begin{center}
        {\includegraphics[height=3cm]{fh-logo}}\\[1cm]
        {\LARGE \thetitle}\\[0.5cm]
        {\large \theauthor}\\
        \href{mailto:97hilfel@gmail.com}{97hilfel@gmail.com}\\
    \end{center}
    \tableofcontents
    \clearpage
}

\newcommand{\screenshot}[1]{
    \begin{figure}[H]
        \centering
        \includegraphics[scale=0.375]{#1}
    \end{figure}
}

% Author and subject
\author{Felix Hillebrand}
\newcommand{\subject}{SSD4UE}

\graphicspath{ {./images/} }

% Main document
\begin{document}

\title{Subject~-~Exercise}
\coverpage

\pagenumbering{roman}
\clearpage
\pagenumbering{arabic}
\pagestyle{main}

\chapter{Topic}
\begin{enumerate}
    % First main item
    \item \label{item:backrefA} \hfill % Use hfill to move the figure to the next line
    \begin{description}
        \item[Aufgabe:] Übersetzen Sie die nachfolgenden textuellen Angaben in entsprechende XQuer\-Ausdrücke,
        sodass die bei jeder Aufgabenstellung angegebene Ausgabe erzeugt wird.
        Testen Sie diese Ausdrücke auf Basis des bereitgestellten XML\-Dokuments ue3\_1\_fitnessdokument.xml.
        Die Verwendung von XQuery 1.0 ist ausreichend.
        Der betroffene Namensraum aus dem XML-Dokument soll jeweils beibehalten werden.

        \begin{itemize}
        \item Selektieren Sie die Messwert-Elemente,
        die zu Gewicht-Messungen gehören und sortieren Sie aufsteigend nach dem Datum der Messung.
        Umgeben Sie das gesamte Resultat mit einem Element <M>.
        \textit{(1,5 Punkt)})

        \inputminted{xquery}{assets/ue3_1_1.xquery}

        \item Selektieren Sie Messung-Elemente mit Blutdruck-Werten und verarbeiten Sie nur Messungen,
        die im Kommentar (Notiz) ‚maschinell‘ (case-insensitiv) enthalten.
        Erstellen Sie ein neues Blutdruck-Element mit dem Zeitpunkt der Messung „zeit“ als Attribut.
        Fügen Sie zwei Kind-Elemente hinzu, „Sys“ für den systolischen Wert und „Dia“ für den diastolischen Wert.
        Packen Sie die Ausgabe in ein „RR“-Element. Verwenden Sie zur Formatierung der Ausgabe Element- und Attributkonstruktoren.
        \textit{(1,5 Punkte)}

        \item Erstellen Sie eine Aktivitätsübersicht für alle Geräte,
        mit denen bereits mindestens 2 Messvorgänge durchgeführt wurden,
        führen Sie dazu auch die Zeitpunkte an.
        Erstellen Sie hierfür ein Element „Aktivität“ und listen Sie darin „Gerät“-Elemente (mit der „id“ als Attribut) auf.
        Die Kindelemente „Online“ von „Gerät“ enthalten die Messzeitpunkte.
        \textit{(1,5 Punkte)}
        \end{itemize}

         \hfill
        \item[Lösung:] \hfill \newline % Use hfill to move the figure to the next line
            Solution
    \end{description}
\end{enumerate}

\begin{enumerate}
    % First main item
    \item \label{item:backrefA} \hfill % Use hfill to move the figure to the next line
    \begin{description}
        \item[Aufgabe:]
         \hfill
        \item[Lösung:] \hfill \newline % Use hfill to move the figure to the next line
            Solution
    \end{description}
\end{enumerate}
\newpage

\chapter{Legal}
Die Ausarbeitung der Aufgabe wurde durch \texttt{OpenAI - GPT-4.5 Turbo}, \texttt{OpenAI - GPT-4.5 Vision}, \texttt{Anthropic -- Claude 3 Opus},  \texttt{Kagi - FastGPT} mit mehreren unterschiedlichen Prompts und Custom Instructions unterstützt.

\end{document}
